\begin{definition} $\PATH = \{(G, s, t) \ | \text{ в орграфе $G$ есть путь из $s$ в $t$}\}$. \end{definition}

\begin{theorem} $\nlsp = \conlsp$ \end{theorem}

\begin{proof} В прошлом пункте была доказана $\nlsp$ полнота $\PATH$, тогда $\overline{\PATH}$ будет $\conlsp$-полным, откуда понятно, что нам хватит доказать, что $\overline{\PATH} \in \nlsp$, тогда любая задача из $\conlsp$ будет сводится к $\nlsp$.

Если известно число всех вершин, достижимых из $s$, то можно показать, что $t$ недостижима. Для этого нам надо предъявить список достижимых из $s$ вершин с их путями. Для этого по индукции будем считать число достижимых за $i$ ходов вершин.

Пусть $c_i$~--- число вершин, достижимых из $s$ за не более, чем $i$ шагов, тогда $c_0 = 1$, а $c_n$~--- размер компоненты достижимости из $s$. Сертификат мы соберем из сертификатов подсчета $c_i$ и того, что $t$ недостижима за $n$ шагов.

База $c_0$ тривиальна. Сделаем переход, пусть мы убедили верификатор в корректности $c_i$. Сертификатом достижимости вершины $v$ будет путь, для его проверки помним $i$, текущую вершину и число ребер в пути. Проверяем для каждого ребра из текущей вершины, что длина пути не превысит $i + 1$, что последняя вершина $v$ и наличие ребра.

С сертификатом недостижимости всё малость труднее. Он состоит из отсортированного списка всех вершин, достижимых за $i$ шагов в соответствующими путями. Для проверки сертификата храним $i$, $c_i$, число уже проверенных вершин, проверяемую вершину, вершину в проверяемом пути до проверяемой вершины и длину проверенной части проверяемого пути, а это все влезает в логарифмическую память.

Проверяем так:
\begin{enumerate}
    \item Новая вершина по номеру больше старой (отсортированность списка)
    \item Проверяем корректность пути из $s$ до проверяемой вершины (последовательное наличие всех ребер и длина проверяемого пути)
    \item Отсутствие ребра из проверяемой вершины в $v$
\end{enumerate}

Тогда верификатор сможет посчитать $c_{i + 1}$ и забыть $c_i$. Корректность следует из конструктивных размышлений.

После такого увлекательного процесса, зная $c_0$, верификатор будет знать $c_{n - 1}$, а потом просто сертификат отсутствия пути до $t$.

Мы смогли показать, что $\overline{\PATH} \in \nlsp$, а значит $\nlsp = \conlsp$. \end{proof}

\begin{statement} $\overline{\PATH} \in \nlcsp$. \end{statement} 

\begin{proof} Из теоремы Иммермана-Селепченьи следует, что $\overline{\PATH} \in \nlsp$, осталось победить полноту. Пусть $A \in \nlsp$, тогда $\overline{A} \in \nlsp$ (по теореме Имммермана-Селепченьи), то есть $\overline{A} \lsv \PATH$ или $A \lsv \overline{\PATH}$, что и требовалось показать. \end{proof}